\section{匹配方案与有限数据预测}\label{sec:implications}

表示结果描述完整的条件通行次序，但经济应用通常只能观察有限多个方案选择。本节直接处理这些观察。主要结果从序数方案数据出发，为反事实通行价值给出锐界；随后用一个简短的货币校准推论说明，同一有限设计几何如何以货币单位表达。

\subsection{来自序数方案选择的锐界}\label{sec:finiteprediction}

固定一个有限物理网格 $\Pi$。令 $x_1,\ldots,x_J$ 表示该网格上的有利单单元通行差分：趋向目标 单元按“通行相对于停留”的方向取向，远离目标 单元按“停留相对于通行”的方向取向。它们张成有限维网格账户空间 $E_\Pi$。令
\[
 e_\Pi=\sum_{j=1}^J x_j.
\]
设 $\mathcal O=\{(\gamma_i,\eta_i)\}_{i=1}^n$ 为网格上已观察的弱方案选择，并假定每个观察对内部维持的背景属性得到控制；记 $d_i=q(\gamma_i)-q(\eta_i)$。定义
\[
 \mathcal P_\Pi(\mathcal O)
 =\left\{
 p\in E_\Pi^*:
 p(e_\Pi)=1,\quad
 p(x_j)\ge0\ \forall j,\quad
 p(d_i)\ge0\ \forall i
 \right\}.
\]
这些弱单元不等式给出有限设计的闭合、面向目标的补全集合。

对同一网格上的一个未观察方案对 $(\gamma,\eta)$，令 $d=q(\gamma)-q(\eta)$，并定义
\[
 \underline\Delta_{\mathcal O}(d)
 =\min_{p\in\mathcal P_\Pi(\mathcal O)}p(d),
 \qquad
 \overline\Delta_{\mathcal O}(d)
 =\max_{p\in\mathcal P_\Pi(\mathcal O)}p(d).
\]

\begin{theorem}[有限设计反事实通行锐界]\label{thm:finite-bounds}
假设 $\mathcal P_\Pi(\mathcal O)\neq\varnothing$。则它是一个紧凸多面体，并且对每个反事实差分 $d$，
\[
 \{p(d):p\in\mathcal P_\Pi(\mathcal O)\}
 =\bigl[\underline\Delta_{\mathcal O}(d),\overline\Delta_{\mathcal O}(d)\bigr].
\]
两个端点分别由两个线性规划问题求得。因此，严格为正的下界稳健地把 $\gamma$ 排在 $\eta$ 之上；严格为负的上界稳健地把 $\eta$ 排在 $\gamma$ 之上；若区间同时包含正负值，则记录了相容可加补全之间的反事实歧义。
\end{theorem}

\begin{proof}
因为 $p(x_j)\ge0$ 且 $p(e_\Pi)=\sum_jp(x_j)=1$，每个单元价值都位于 $[0,1]$。又由于 $x_j$ 张成 $E_\Pi$，归一化泛函有界；其余限制只是有限多个闭线性不等式。因此 $\mathcal P_\Pi(\mathcal O)$ 是紧凸多面体。在映射 $p\mapsto p(d)$ 下，其线性像为一个紧区间，端点由上述两个线性规划达到。多面体中的每个元素都定义有限观察的一个弱面向目标可加补全，由此得到锐性。
\end{proof}

定理~\ref{thm:D} 是精确转移的特例：当一个已观察比较对与一个反事实比较对具有相同的通行账户差分时，同一可加限制同时适用于二者。当不存在任何一个单独的已观察比较拥有目标差分时，有限数据定理还可以组合多个绕行、反向和目标接触选择的信息。

\subsection{货币校准}\label{sec:wtp}

有限的“方案--价格”数据可以在不为完整原始次序选定单一支撑赋值的情况下，把同样的可加对比校准为货币单位。令一个观察对象为 $(\gamma,m)$，并在该有限设计内维持货币拟线性。相容的可加通行赋值满足
\[
 p\bigl(q(\gamma_i)-q(\eta_i)\bigr)\ge m_i-n_i
\]
对每一个已观察弱选择 $(\gamma_i,m_i)\succeq(\eta_i,n_i)$ 成立，并同时满足所选网格上的目标方向不等式。

\begin{corollary}[货币校准的通行锐界]\label{cor:wtp}
对任意反事实通行差分 $d$，在相容有限设计多面体上，货币校准值集合 $\{p(d)\}$ 是一个区间，可能无界；其有限端点由两个线性规划求得。若相对于 $\eta$，方案 $\gamma$ 还需要支付额外费用 $c$，则严格高于 $c$ 的下界给出对 $\gamma$ 的稳健严格偏好，而严格低于 $c$ 的上界给出对 $\eta$ 的稳健严格偏好。
\end{corollary}

该推论只是对有限设计识别集合的货币校准。本文的经济比较并不依赖全局拟线性补全，方案--价格计算也没有被用于完整定义域的表示定理。

\subsection{识别与适用范围}

固定有限物理网格会消除非平凡的穿孔收缩序列，因此 PROC 正则性问题在该网格上成为空洞条件。当不受限路径域允许无界重复时，PVBC 仍有实质内容：注~\ref{rem:finitegrid-pvbc} 给出了一个满足有限消去与目标方向、但违反 PVBC 的字典序有限网格次序。相比之下，有限经验设计只使用有限多个已观察不等式，其可加补全问题按构造就是多面体问题。

在序数方案数据下，与决策相关的识别对象通常是
\[
 [\underline\Delta_{\mathcal O}(d),\overline\Delta_{\mathcal O}(d)].
\]
在方案--价格数据下，同一有限设计不等式把这个对比区间校准到货币单位。无论是序数区间还是货币校准区间，都可以在不分别识别 $\alpha_s$、$\kappa_s$ 或 $P$ 的情况下直接用于决策。精确无差异可以使序数相容多面体收缩。在一个 $d$ 维设计上，$d-1$ 个独立无差异等式加上归一化足以固定唯一可行的可加补全；更一般地，单点识别应根据全部有效约束系统检查。

条件通行次序最直接地由受控方案比较观察得到。当维持的背景属性被匹配或单独净除时，序数转移与归一化有限网格界均适用。方案--价格菜单提供一种明确方式来放松“货币成本相同”的要求：费用通过拟线性计价物进入，而通行分量仍保持集合识别。其他健康、持续时间和动态后果，在其价值已知或已建模时，可以通过应用特定结构加入。

BPCC 仍是不受限路径域上的一个全局有限列表限制。附录~\ref{app:primitive} 给出一个六历史族，其中原始平衡证书需要无界的重复次数。在有限经验设计中，可加可行性则直接通过相应的线性不等式系统检验。
